\begin{remark}
\label{rem:wt-unwt}
  It has been shown by Crescenzi, Silvestri and Trevisan~\cite{CST01}
  that any hardness result for weighted instances of \maxcut{k}
  carries over to  unweighted instances {assuming} the
  total edge weight is polynomially bounded. In fact, their reduction
  preserves $k$-colorability, so an inapproximability result for the weighted
  \maxcolor{k} problem also holds for the unweighted version. Therefore
  all our hardness results hold for the unweighted \maxcolor{k} problem.
\end{remark}

\section{Conditional hardness results for \maxcolor{k}}
%
We will first review the \dtoone{2} Conjecture, and then construct a
noise operator, which allows us to preserve $k$-colorability. Then we
will bound the stability of coloring functions with respect to this
noise operator. In the last section, we will give a PCP verifier which
concludes the hardness result.
%
\vspace{-1ex}
\subsection{Preliminaries}
\label{sec:condhardness}
%
We begin by reviewing some definitions
%
and the \dtoone{d} conjecture.
%
\begin{definition}
  An instance of a bipartite Label Cover problem represented as
  $\mathcal{L} = (U, V, E, W, R_U, R_V, \Pi)$ consists of a weighted
  bipartite graph over node sets $U$ and $V$ with edges $e=(u,v) \in
  E$ of non-negative real weight $w_e \in W$. $R_U$ and $R_V$ are
  integers with $1\le R_U\le R_V$. $\Pi$ is a collection of projection
  functions for each edge:
  \[
  \Pi=\{\pi_{vu}:\intset{R_V}\rightarrow\intset{R_U}\;\big|\;u\in U, v\in
  V\}\,.
  \]
  A labeling $\ell$ is a mapping $\ell:U\rightarrow
  \intset{R_U}$, $\ell:V\rightarrow \intset{R_V}$. An edge $e=(u,v)$
  is satisfied by the labeling $\ell$ if $\pi_e(\ell(v))=\ell(u)$. We
  define the value of a labeling as the sum of weights of edges
  satisfied by this labeling normalized by the total
  weight. $\opt(\mathcal{L})$ is the maximum value over any labeling.
\end{definition}
%
%
\begin{definition}
  A projection $\pi:\intset{R_V}\rightarrow \intset{R_U}$ is called
  \dtoone{d} if for each $i \in \intset{R_U}$, $|\pi^{-1}(i)|\le
  d$. It is called \emph{exactly} \dtoone{d} if $|\pi^{-1}(i)|= d$ for
  each $i \in \{1,2,\dots,R_U\}$.
\end{definition}
%
%
\begin{definition}
  A bipartite Label-Cover instance $\mathcal{L}$ is called $d$-to-$1$
  Label-Cover if all projection functions, $\pi \in \Pi$ are
  $d$-to-$1$.
\end{definition}
%
The following conjecture was first introduced in a seminal paper by
Khot~\cite{Khot02}. We use a slight refinement of it due to
Dinur et al.~\cite{DMR}.
\begin{conjecture}[\dtoone{d} Conjecture, as stated in~\cite{DMR}]
  For every constant $\gamma > 0$, the following decision problem is \cclass{NP}-hard.
  Given an exactly \dtoone{d} Label-Cover instance
  $\mathcal{L}$ with $R_V=R(\gamma)$ and $R_U = d R_V$ many labels:
  \begin{itemize}
  \item {\bf Yes:} $\opt(\mathcal{L}) = 1$,
  \item {\bf No:} $\opt(\mathcal{L})\le \gamma$.
  \end{itemize}
%
%
%
%
%
%
\label{conj:dtoone}
\end{conjecture}
%
Using the reductions from~\cite{DMR}, it is possible to show that the
above conjecture still holds given that the graph $(U\cup V, E)$ is
left-regular and unweighted, \ie, $w_e = 1$ for all $e\in E$.
%
\subsection{Noise  operators}
%
%
%
%
%
For a positive integer $M$, we will denote by $[M]$ the set
$\{0,1,\dots,M-1\}$.  We will identify elements of $[M^2]$ with $[M]
\times [M]$ in the obvious way, with the pair $(a,b) \in [M]^2$
corresponding $a + M b \in [M^2]$. 
%
%
\begin{definition}
  \label{def:markovop}
  A Markov operator $T$ is a linear operator which maps probability
  measures to other probability measures. In a finite discrete
  setting, it is defined by a stochastic matrix whose $(x,y)$'th entry
  $T(x \rightarrow y)$ is the probability of transitioning from $x$ to
  $y$. Such an operator is called symmetric if $T(x \rightarrow y) =
  T(y \rightarrow x) = T(\xiffy{x}{y})$.
\end{definition}
\begin{definition}
  \label{def:beckner}
  Given $\rho\in [-1,1]$, the Beckner noise operator, $T_{\rho}$ on
  $[q]$ is defined as \[
  T_\rho(x\rightarrow x) = \frac{1}{q} +
  \left(1-\frac{1}{q}\right) \rho\quad \text{and}\quad T_\rho(x\rightarrow y) =
  \frac{1}{q}(1-\rho)\] for any $x\neq y$. 
\end{definition}
%
\begin{observation}\label{obs:basis-id}
  All eigenvalues of the operator $T_{\rho}$ are given by \[1 =
  \lambda_0(T_\rho)\ge \lambda_1(T_\rho) = \cdots =
  \lambda_{q-1}(T_\rho) = \rho\,.\] Any orthonormal basis
  $\alpha_0,\alpha_1,\ldots,\alpha_{q-1}$ with $\alpha_0$ being
  constant vector, is also a basis for $T_{\rho}$.
\end{observation}
%
\begin{lemma}
  \label{lemma:T}
  For an integer $q \ge 6$, there exists a symmetric Markov operator
  $T$ on $[q]^2$ whose diagonal entries are all $0$ and with
  eigenvalues $1=\lambda_0\ge \lambda_1\ge \ldots\ge \lambda_{q^2-1}$
  such that the \emph{spectral radius} $\rho(T) =
  \max\{|\lambda_1|,|\lambda_{q^2-1}|\}$ is at most
  ${4}/{(q-1)}$.%
\end{lemma}
\begin{proof}
  Consider the symmetric Markov operator $T$ on $[q]^2$ such that, for
  $x=(x_1,x_2), y=(y_1,y_2)\in [q]^2$,
  \[T(\xiffy{x}{y}) =
  \begin{cases}
    \alpha &\mbox{if $\{x_1,x_2\}\cap\{y_1,y_2\}=\emptyset$
      and $x_1\neq x_2, y_1\neq y_2$,}\\
    \beta &\mbox{if $x_1\not\in\{y_1,y_2\}$
      and $x_1= x_2, y_1\neq y_2$,}\\
    \beta &\mbox{if $y_1\not\in\{x_1,x_2\}$
      and $x_1\neq x_2, y_1= y_2$,}\\
    %
    %
    0 & \mbox{else,}
  \end{cases}\] where \[\alpha = \frac{1}{(q-1)(q-3)}\quad \text{and}\quad \beta =
  \frac{1}{(q-1)(q-2)}\,.\] It is clear that $T$ is symmetric and doubly
  stochastic.

  To bound the spectral radius of $T$, we will bound the second
  largest eigenvalue $\lambda_1(T^2)$ of $T^2$. Notice that $T^2$ is
  also a symmetric Markov operator. Moreover the set of eigenvalues of
  $T^2$, $\{\lambda_i(T^2)\}_i$ is equal to $\{{\lambda_i(T)}^2\}_i$.
  %
  Therefore \[\lambda_1(T^2) \ge
  \max(\lambda^2_1(T),\lambda^2_{q^2-1}(T)) \ge \rho(T)^2\,.\]

  Notice that $T^2(\xiffy{x}{y})>0$ for all pairs $x,y\in[q]^2$.
  %
  %
  %
  %
  Consider the variational characterization of $1 - \lambda_1(T^2)$
  \cite{Sinclair92}:
  \[\min_\psi \frac{\sum_{x,y}
    (\psi(x) - \psi(y))^2 \pi(x) T^2(\xiffy{x}{y})}{\sum_{x,y}
    (\psi(x)-\psi(y))^2 \pi(x)\pi(y)} \ge \min_\psi \min_{x,y}
  \frac{\pi(x) (\psi(x) - \psi(y))^2
    T^2(\xiffy{x}{y})}{(\psi(x)-\psi(y))^2 \pi(x)\pi(y)} = \min_{x,y}
  q^2 T^2(\xiffy{x}{y})\] where $\pi$ corresponds to the stationary
  distribution for Markov operator $T^2$.

  %

  For any two pairs $(x_1,x_2),(y_1,y_2)\in[q]^2$, let $l =
  |[q]\setminus \{x_1,x_2,y_1,y_2\}|$. Then we have
  \begin{eqnarray*}
    T^2(\xiffy{(x_1,x_2)}{(y_1,y_2)}) &=& 
    \begin{cases} l (l-1)\beta^2\ge (q-2)(q-3)\beta^2 & 
      \mbox{if $x_1=x_2$ and $y_1=y_2$,}\\
      l (l-1) \alpha \beta\ge (q-3)(q-4)\alpha\beta & 
      \mbox{if $x_1\neq x_2$ and $y_1=y_2$,}\\
      l (l-1) \alpha \beta\ge (q-3)(q-4)\alpha\beta & 
      \mbox{if $x_1=x_2$ and $y_1\neq y_2$,}\\
      l (l-1) \alpha^2 + l \beta^2 \ge (q-4)((q-5)\alpha^2 + \beta^2) &
      \mbox{if $x_1\neq x_2$ and $y_1\neq y_2$.}
  \end{cases}\\
  &\ge& \frac{(q-5) (q-4)}{(q-3)^2 (q-2) (q-1)}\,.
  \end{eqnarray*}
  So \[\rho(T)\le \sqrt{\lambda_1(T^2)} \le \sqrt{1-\frac{(q-5) (q-4)
      q^2}{(q-3)^2 (q-2) (q-1)}} \le 
%
\frac{4}{q-1}\quad \text{for} \quad q\ge 6\,.\qedhere\]
\end{proof}
%
\subsection{\texorpdfstring{$q$-ary functions, influences, noise stability}{q-ary functions, influences, noise stability}}
%
We define inner product on space of functions from $[q]^N$ to
$\R$ as $\langle f,g \rangle = \expct{x\sim[q]^N}{f(x)
  g(x)}$. Here $x \sim \mathcal{D}$ denotes sampling from distribution
$\mathcal{D}$ and $\mathcal{D} = [q]^N$ denotes the uniform
distribution on $[q]^N$.

Given a symmetric Markov operator $T$ and
$x=(x_1,\ldots,x_N)\in[q]^N$, let $T^{\otimes N} x$ denote the product
distribution on $[q]^N$ whose \texsup{$i$}{th} entry $y_i$ is
distributed according to $T(\xiffy{x_i}{y_i})$. Therefore $T^{\otimes
  N} f (x) = \expct{y \sim T^{\otimes N} x} { f(y)}$.

\begin{definition}
  Given $x\in [q]^N$ and $a\in[q]$, let $|x|_a$ be the number of coordinates
equal to $a$, $a=\left| \left\{ i\mid x_i = a\right\}\right|$.
We will use $|x|$ to denote the number of non-zero coordinates of $x$,
$|x| = \sum_{a\neq 0} |x|_a$.
\end{definition}
\begin{definition}
  Let $\alpha_0,\alpha_1,\ldots,\alpha_{q-1}$ be an orthonormal basis
  of $\R^q$ such that $\alpha_0$ is all constant vector. For
  $x\in[q]^N$, we define $\alpha_x \in \R^{q^N}$ as
  \[ \alpha_x = \alpha_{x_1}\otimes \ldots \otimes \alpha_{x_N}.\] 
\end{definition}
\begin{definition}[Fourier coefficients]
  For a function $f:[q]^N\to\R$, define $\hat{f}(\alpha_x) =
  \langle f, \alpha_x \rangle$.
\end{definition}
%
\begin{definition}
  Let $f:[q]^N\rightarrow \R$ be a function. The
  \emph{influence} of \texsup{i}{th} variable on $f$, $\infl_i(f)$ is
  defined by
  \[ \infl_i(f) = 
      \expct{}
            {\var{}{f(x)|x_1,\ldots,x_{i-1},x_{i+1},\ldots,x_N}}\]
  where $x_1,\ldots,x_N$ are uniformly distributed.
%
\end{definition}
\begin{observation}
  $\infl_i (f) = \sum_{x:x_i\neq 0} \hat{f}^2(\alpha_x)$.
\end{observation}
%
\begin{definition}
  Let $f:[q]^N\rightarrow \R$ be a function. The \emph{low-level
    influence} of \texsup{i}{th} variable of $f$ is defined by
  \[ \infl_i^{\le t}(f) = \sum_{x:x_i\neq 0,\ |x|\le t} \hat{f}^2(\alpha_x).\]
\end{definition}
%
\begin{observation}
  For every function $f:[q]^N\to\mathbb{R}$, \[\sum_i \infl_i^{\le
    t}(f) = \sum_{x:|x|\le t} \hat{f}^2(\alpha_x) |x|\le t \sum_x
  \hat{f}^2(\alpha_x) = t \|f\|^2_2.\] If $f:[q]^N\to[0,1]$, then
  $\|f\|^2_2 \le 1$, so $\sum_i \infl_i^{\le t}(f) \le t$.
\end{observation}
%
\begin{definition}[Noise stability]
  Let $f$ be a function from $[q]^N$ to $\R$, and let $-1\le \rho\le
  1$. Define the \emph{noise stability} of $f$ at $\rho$ as
  \[ \mathbb{S}_{\rho}(f) = \langle f, T^{\otimes n}_\rho f\rangle =
  \sum_x \rho^{|x|} \hat{f}^2_i(\alpha_x)\] where $T_{\rho}$ is the
  Beckner operator as in \expref{Definition}{def:beckner}.
\end{definition}
%
A natural way to think about a $q$-coloring function is as a
collection of $q$-indicator variables summing to $1$ at every point.
To make this formal:
%
\begin{definition}
  Define the unit $q$-simplex as $\Delta_q = \{(x_1,\ldots,x_{q}) \in
  \R^q \mid \sum x_i = 1, x_i \ge 0\}$.
\end{definition}
%
\begin{observation}\label{obs:infbndsimplex}
  For positive integers $Q,q$ and any function
  $f=(f_1,\ldots,f_q):[Q]^N\to \Delta_q$, \[\sum_i \infl^{\le t}_i(f) =
  \sum_i \sum_j \infl^{\le t}_i (f_j) \le t \sum_j \|f_j\|^2 \le t\,.\]
\end{observation}
%
%
We want to prove a lower bound on the stability of $q$-ary functions
with noise operators $T$. The following proposition is generalization
of Proposition 11.4 in~\cite{KKMO} to general symmetric Markov
operators $T$ with small spectral radii.
%
%
\begin{definition}
  \label{def:gamma-rho}
  Let $\rho\in[0,1]$ and $\mu\in[0,1]$. Let $X$ and $Y$ denote
normal random variables with mean $0$ and covariance matrix 
\[
\begin{pmatrix}
  1 & \rho \\ \rho & 1
\end{pmatrix}\,.
\]
We define
\[
\Lambda_{\rho}(\mu) = \mathrm{Prob}\left[X\ge t\mbox{ and }Y\ge
  t\right].
\]
where $t$ is chosen so that $\mathrm{Prob}[X\ge t] = \mu$.
\end{definition}
Following is a restatement of the famous Majority is stablest theorem:
\begin{theorem}[Majority is Stablest, \cite{KKMO}]
  Fix $q\ge 2$ and $\rho\in[0,1)$. Then for any $\eps>0$ there exists
  a small enough $\delta=\delta(\eps,p,q)>0$ such that if
  $f:[q]^n\to[0,1]$ is any function satisfying $\mathbb{E}[f] = \mu$
  and $\infl_i(f)\le \delta$ for all $i=1,\ldots,n$ then
\[
\mathbb{S}_{\rho}(f)\le \Lambda_{\rho}(\mu)+\eps.
\]
\end{theorem}
\begin{proposition}
\label{prop:stab}
For integers $Q,q \ge 3$, and a symmetric Markov operator $T$ on $[Q]$
with spectral radius $\rho(T)\le {c}/{(q-1)}$, for some $c>0$, there
exists $t=t(q)>0$ and a small enough $\delta=\delta(q)>0$ such
that %
any function $f=(f_1,\ldots,f_q):[Q]^N\rightarrow \Delta_q$ with
$\infl_i^{\le t} (f)\le \delta$, for all $i$, satisfies
  \[ \sum_{j=1}^q 
\langle f_j, T^{\otimes N} f_j\rangle 
%
\ge 1/q - 2 c \ln q/q^2 - C
  \ln\ln q /q^2\] for some universal constant $C<\infty$.
\end{proposition}
\begin{proof}
  Let $t=4$, $f_i:[Q]^N\rightarrow[0,1]$ denote the \texsup{$i$}{th}
  coordinate function of $f$, and let $\mu_i = \expct{}{f_i}$. Let
  $\alpha_0,\ldots, \alpha_{Q-1}$ be an orthonormal set of
  eigenvectors for $T$ with corresponding eigenvalues ${\lambda_0\ge
    \ldots\ge \lambda_{Q-1}}$, with $\rho=\rho(T)\le {c}/{(q-1)}$
  being the spectral radius of $T$. Notice that $T$ is symmetric so
  $\lambda_0 = 1$ and $\alpha_0$ is a constant vector. Therefore
  $\expct{}{f_i} = \hat{f}_i(\alpha_0) = \mu_i$. Then (using the
  notation from~\cite{DMR}):
  \[ T^{\otimes N} \alpha_x = (\prod_{a \neq 0} \lambda_a ^{|x|_a}) \alpha_x \]
  and hence
  \[ T^{\otimes N} f_i = \sum_x (\prod_{a \neq 0} \lambda_a ^{|x|_a})
  \hat{f}_i(\alpha_x) \alpha_x. \]

  %
  %
  %
  %
  %
  %

  At this point, consider the Beckner operator, $T_{\rho}$ on
  $[Q]$. Since $\alpha_0$ is the uniform distribution 
  %
  $\alpha_0,\alpha_1,\ldots,\alpha_{Q-1}$ is
  also an orthonormal basis for $T_{\rho}$ by \expref{Observation}{obs:basis-id}. 
  %
  %
  %
  %
  %
  %
  %
  %
  %
  %
  Thus
  %
  \begin{eqnarray*}
    \langle f_i, T^{\otimes N} f_i\rangle &=& 
    \hat{f}_i^2(\alpha_0) - \hat{f}_i^2(\alpha_0) + 
    \sum_x \underbrace{(\prod_{a \neq 0} \lambda_a ^{|x|_a})}_
    {
      \begin{cases}
        \ge - \rho^{|x|} & \text{if }|x|\neq 0,  \\
        = 1 & \text{else.}
      \end{cases}
    }
    \hat{f}_i^2(\alpha_x)\\
    &\ge& 2 \mu_i^2 - 
    \sum_{x} \rho^{|x|} \hat{f}_i^2(\alpha_x) =
    2\mu_i^2 -\sum_{x:|x|\le 4} \rho^{|x|} \hat{f}_i^2(\alpha_x) 
    -\sum_{x:|x|> 4} \rho^{|x|} \hat{f}_i^2(\alpha_x)\\
    &\ge& 2\mu_i^2 
    -\sum_{x:|x|\le 4} \rho^{|x|} \hat{f}_i^2(\alpha_x) 
    - \rho^{4}\\
    &\ge&  2\mu_i^2 
    -\sum_{x:|x|\le 4} \rho^{|x|} \hat{f}_i^2(\alpha_x) 
    - q^{-3}\,.\\
  \end{eqnarray*}
  At this point, let $\tilde{f}_i(x) = \sum_{x:|x|\le 4} (\prod_{a
    \neq 0} \lambda_a ^{|x|_a}) \hat{f}_i(\alpha_x)\alpha_x $ be the
  function having the same low-level coefficients with $f_i(x)$ and
  $0$ for the higher-levels. It is easy to verify that
  $\expct{}{\tilde{f}_i} = \mu_i$, $\infl_i(f_j) \ge
  \infl_i(\tilde{f}_j) = \infl_i^{\le 4}(f_j)$ and
  $\mathbb{S}_{\rho}(\tilde{f_j}) = \sum_{x:|x|\le 4} \rho^{|x|}
  \hat{f}_i^2(\alpha_x)$. In particular, our assumption $\sum_j
  \infl_i^{\le t}(f_j)=\sum_j \infl_i^{\le 4}(f_j) \le \delta$ implies
  $\sum_j \infl_i(\tilde{f}_j) \le \delta$.

  Let $\delta$ be a small enough constant such that
  $\mathbb{S}_{\frac{c}{q-1}}(\tilde{f}_i) \le
  \Lambda_{\frac{c}{q-1}}(\mu_i)+\eps$ for some small $\eps\le
  {1}/{q^3}$, from the Majority is Stablest Theorem~\cite{MOO}.
  %
  %
  Below, for a real $x$, $[x]^+$ denotes $\max \{ x, 0\}$. Then
  \begin{eqnarray*}
    \sum_i \langle f_i, T^{\otimes N} f_i\rangle &\ge& \sum_i \left[2 \mu_i^2
      - \mathbb{S}_{\rho}(\tilde{f}_i)\right]  - q^{-2}\\
    &\ge& \sum_i \left[2 \mu_i^2 - \mathbb{S}_{\frac{c}{q-1}}(\tilde{f}_i)\right]  -
    q^{-2}\\
    &\ge& \sum_i \left[2 \mu_i^2 - \Lambda_{\frac{c}{q-1}}(\mu_i)\right]^+-2 q^{-2}\\
    &\ge& \frac{1}{q} - \frac{2 c \ln q}{q^2} - O\left(\frac{\ln\ln q}{q^2}\right)\,.
  \end{eqnarray*}
  The last inequality is proved in the same way as Proposition 11.4 in~\cite{KKMO}. The only difference is that we have
  \[ F(\mu_i) = \mu_i^2 + \frac{c}{q-1} 2 \mu_i^2 \ln(1/\mu_i)\cdot
  \left(1+C\frac{\ln\ln q}{\ln q}\right)\]
  and
  \[ \sum_{i=1}^q \left[ 2 \mu_i^2 -
    \Lambda_{\frac{c}{q-1}}(\mu_i)\right]^+ \ge \sum_{i=1}^q (2\mu_i^2
  - F(\mu_i))\] which is convex because $\mu_i\le (1/q)^{1/10}$ and
  minimized at $\mu_i=1/q$. In this case, we have
  \[ \sum_{i=1}^q (2\mu_i^2 - F(\mu_i)) \ge q\left(q^{-2} - 2 c q^{-3}
    \ln q (1 + C \ln \ln q / \ln q\right)\] from which the above claim
  follows.
\end{proof}

%
\begin{definition}[Moving between domains]
  For every $x=(x_1,\ldots,x_{2 N}) \in [q]^{2 N}$, denote
  $\overline{x} \in [q^2]^N$ as
\[ \overline{x} = ((x_1, x_2),\ldots,(x_{2 N-1},x_{2 N})) \ . \]
  Similarly for $y=(y_1,\ldots,y_{N}) \in [q^2]^{N}$, denote
  $\underline{y} \in [q]^{2 N}$ as
  \[ \underline{y} = (y_{1,1}, y_{1,2},\ldots,y_{N,1},y_{N,2})\,,\]
  where $y_i = y_{i,1} + y_{i,2} q$ such that $y_{i,1},y_{i,2}\in
  [q]$.  For a function $f$ on $[q]^{2 N}$, define $\overline{f}$ on
  $[q^2]^N$ as $\overline{f}(y) = f(\underline{y})$. 
%
  %
\end{definition}
%
\noindent 
The relationship between influences of variables for functions $f$ and
$\overline{f}$ are given by the following claim (Claim 2.7 in~\cite{DMR}).
%
\begin{claim}\label{claim:infrel}
  For every function $f:[q]^{2N}\rightarrow \R$,
  $i\in\intset{N}$ and any $t\ge 1$,
  \[\infl_i^{\le t}(\overline{f}) \le \infl_{2i-1}^{\le 2t}(f) +
  \infl_{2i}^{\le 2t} (f)\,.\]
\end{claim}

%
%
%
%
%
%
%
%
%
%
%
%
%
%
%
%
%
%
%
%
%
%
%
%
%
%
%
%
%
%
%
%
%
%
%
%
%
%
%
%
%


\vspace{-1ex}
\subsection{PCP verifier for \maxcolor{k}}
%
This verifier uses ideas similar to the \maxcut{k} verifier given in~\cite{KKMO} and the $4$-coloring hardness reduction in~\cite{DMR}. Let $\mathcal{L} = (U,V,E,R,2 R,\Pi)$ be a \dtoone{2}
bipartite, unweighted and left regular Label-Cover instance as in
\expref{Conjecture}{conj:dtoone}. Assume the proof is given as the Long
Code over $[k]^{2 R}$ of the label of every vertex $v \in V$. Below
for a permutation $\sigma$ on $\intset{n}$ and a vector $x \in \R^n$, $x\circ \sigma$ denotes $(x_{\sigma(1)},x_{\sigma(2)},\cdots,x_{\sigma(n)})$. For a function $f$ on $\R^n$, $f \circ \sigma$ is defined as $f \circ \sigma(x) = f(x \circ \sigma)$.

%
%

%
%

\begin{itemize}
\vspace{-1ex}
\itemsep=0ex
\item Pick $u$ uniformly at random from $U$, $u\sim U$.
\item Pick $v,v'$ uniformly at random from $u$'s neighbors. Let $\pi,
  \pi'$ be the associated projection functions, $\chi_v, \chi_{v'}$ be
  the (supposed) Long Codes for the labels of $v, v'$ respectively.
\item Let $T$ be the Markov operator on $[k]^2$ given in \expref{Lemma}{lemma:T}. Pick $x \sim [k^2]^R$ and $y \sim T^{\otimes R}
  x$. Let $\sigma_v,\sigma_{v'}$ be two permutations of $\intset{2 R}$ such
  that \[\pi(\sigma^{-1}_v(2i-1)) = \pi(\sigma^{-1}_v(2i)) =
  \pi'(\sigma_{v'}^{-1}(2i-1)) = \pi'(\sigma_{v'}^{-1}(2i))\] (both $\pi$ and
  $\pi'$ are exactly \dtoone{2}, so such permutations exist).
\item Accept iff $\chi_v\ \circ\ \sigma_v (\underline{x}) $ and
  $\chi_{v'} \ \circ\ \sigma_{v'}(\underline{y})$ are different.
\end{itemize}
%
%
%
\begin{lemma}[Completeness]
\label{lem:kcol-compl}
  If the original \dtoone{2} Label-Cover instance $\mathcal{L}$ has a
  labeling which satisfies all constraints, then there is a proof
  which makes the above verifier always accept.
\end{lemma}
%
\begin{proof}
  Let $\ell:V\to \intset{2 R}$ be a labeling for $\mathcal{L}$
  satisfying all constraints in $\Pi$. Pick $\chi_v$ as the Long Code
  encoding of $\ell(v)$. Given any pair of vertices $v,v' \in V$ which
  share a common neighbor $u \in U$, and $x,y\in [k]^{2 R}$ pairs such
  that 
\[ \prob{}{\overline{y} \sim T^{\otimes R}(\overline{x})} = \prod_i
  T(\xiffy{(x_{2i-1},x_{2i})}{(y_{2i-1},y_{2i})})>0 \ , \]
%
 let $\pi,\pi'$
  be the projection functions and $\sigma_v,\sigma_{v'}$ be the permutations
  as defined in the description of the verifier. We have $\chi_v(x \ \circ\ \sigma_v) =
  x_{\sigma(\ell(v))}$ and $\chi_{v'}(y \ \circ\ \sigma_{v'}) =
  y_{\sigma'(\ell(v'))}$. Since $\pi(\ell(v)) = \pi'(\ell(v'))$, this
  implies $\sigma_v(\ell(v)),\sigma_{v'}(\ell(v')) \in \{2i-1,2i\}$ for some
  $i\le R$. But 
\[ T(\xiffy{(x_{2i-1},x_{2i})}{(y_{2i-1},y_{2i})})>0
  \implies \{x_{2i-1},x_{2i}\}\cap \{y_{2i-1},y_{2i}\}=\emptyset \ , \]
%
  therefore $\chi_{v}\circ \sigma_v(x) = x_{\sigma_v(\ell(v))} \neq
  y_{\sigma_{v'}(\ell(v'))}=\chi_{v'}\circ \sigma_{v'}(y)$. So the verifier always accepts.
\end{proof}
%
\begin{lemma}[Soundness]
\label{lem:kcol-soundness}
  There is a constant $C$ such that, if the above verifier passes with
  probability exceeding $1-1/k+O(\ln k/k^2)$, then there
  is a labeling of $\mathcal{L}$ which satisfies $\gamma'=\gamma'(k)$
  fraction of the constraints independent of label set size $R$.
\end{lemma}
%
\begin{proof}
  For each node $v \in V$, let $f^v:[k]^{2 R} \rightarrow \Delta_k$ be
  the function $f^v(x) = e_{\chi_v(x)}$ where $e_i$
%
is the indicator
  vector of the \texsup{$i$}{th} coordinate. Let $\Gamma(u)$ denote
  the set of vertices adjacent to $u$ in the Label Cover graph.

  After arithmetizing, we can write the verifier's acceptance probability as
%
%
  \[
    \begin{array}{rclr}
      \prob{}{\text{acc}} &=&
      \expct{u,v,v'} {1-
        \sum_j \langle \overline{f^v_j \circ\ \sigma_v},
        T^{\otimes R} \overline{(f^{v'}_j \circ\ \sigma_{v'})}\rangle} \\
      &=& 1 - \expct{u}{\sum_j \expct{v,v'} {
          \langle \overline{f^v_j\ \circ\ \sigma_v},
          T^{\otimes R} \overline{(f^{v'}_j\ \circ\ \sigma_{v'})}\rangle}} \\
      &=& 1 - \expct{u} {\sum_j
        \langle \expct{v}{\overline{f^v_j\ \circ\ \sigma_v}},
        T^{\otimes R} \expct{v'}{
          \overline{f^{v'}_j\ \circ\ \sigma_{v'}}}\rangle} \\
      &=& 1 - \expct{u} {
        \sum_j\langle g^u_j, T^{\otimes R} g^u_j\rangle} & 
      \left(g^u_j = \expct{v\sim \Gamma(u)}
        {\overline{f^v_j\circ \sigma_v}}\right)\\
      &\ge& 1 - 1/k + C \ln k/k^2 
    \end{array}
    \] where $g^u:[k^2]^{R}\to \Delta_{k}$ and some constant $C$.
    %
    By averaging, for at least a fraction $\delta=(\eps/2) \ln k/k^2$
    of vertices in $U$, we have
    \[
    \sum_j\langle g^u_j, T^{\otimes R} g^u_j\rangle \le 1/k-C \ln
    k/k^2\,.\]
    %
    Let these be ``good'' vertices. For a good vertex, by \expref{Proposition}{prop:stab}, there exist constants $\delta=\delta(k)$,
    $t=t(k)$ and $i$ such that $\infl_i^{\le t}(g^u)\ge\delta$. Let
    \[\sugg_u = \{i\mid i\in \intset{R}\wedge \infl_{i}^{\le t} (g^u)\ge
    \delta\}\,,\] so $|\sugg_u|\ge 1$. By \expref{Observation}{obs:infbndsimplex}, $|\sugg_u|\le t/\delta$. For a good
    vertex $u$, and $j\in\sugg_u$:
    %
    \[ \delta \le \infl_j^{\le t} (g^u) = \expct{v \sim
      \Gamma(u)}{\infl_{j}^{\le t} \bigl( \overline{f^v \circ \sigma_v} \bigr)} \] 
    %
    Therefore, for at least a fraction $\delta/2$ of neighbors $v$ of
    $u$, $\infl^{\le t}_{j}(\overline{f^v \circ \sigma_v}) \ge
    \delta/2$.  For such $v$ and $j$, by \expref{Claim}{claim:infrel},
    \[
\infl_{2j-1}^{\le 2t} (f^v\circ \sigma_v) + \infl_{2j}^{\le 2t}
    (f^v\circ \sigma_v) \ge \delta/2\,.
\] Therefore for some $j\in[2 R]$,
    $\infl_j^{\le 2t} (f^v) \ge \delta/4$. Let \[\sugg_v = \{j\mid j\in
    \intset{2 R}\wedge \infl_{j}^{\le 2t} (f^v)\ge \delta/4\}\,.\] Again,
    $\sugg_v$ is not empty and $|\sugg_v|\le 8 t/\delta$.

    Following the decoding procedure in~\cite{KKMO}, we deduce that it
    is possible to satisfy a fraction $\gamma' = \gamma'(\delta,t) =
    \gamma'(k)$
    %
    of the constraints.
\end{proof}

%
Note that our PCP verifier makes ``$k$-coloring'' tests. By the
standard conversion from PCP verifiers to CSP hardness, and
\expref{Remark}{rem:wt-unwt} about conversion to unweighted graphs with
the same inapproximability factor, we conclude the main result of this
section by combining Lemmas \ref{lem:kcol-compl} and
\ref{lem:kcol-soundness}.
\begin{theorem}
  For every constant $k \ge 3$, assuming the \dtoone{2} Conjecture, it
  is \cclass{NP}-hard to approximate \maxcolor{k} within a factor of
  $1-1/k+O(\ln k/k^2)$.
  %
  %
  %
\end{theorem}

